\section{The growth of classes admitting an improper symmetry}
In even positive rank $n\ge4$ let $b^-(m)$ count the ordinary represented classes whose intrinsic automorphism group contains an element of determinant $-1$; no requirement is made that it fix the parent. Then for $X\ge2$, as $X$ tends to infinity,
\begin{equation}\label{eq:reflection-growth}
 X^{n-1}\ll_L\sum_{m\le X}b^-(m)\ll_{L,\epsilon}X^{n-1+\epsilon}.
\end{equation}
The series with coefficients $b^-(m)$ has abscissa $n-1$ and diverges there. Thus the absolute-convergence half-plane for the improper term of Corollary~\ref{cor:KLL} is sharp.

The coefficient of the improper transporter sum is $b^-(m)/2$,
since an intrinsic automorphism group containing an improper element
has equally many automorphisms of each determinant sign. The upper
bound follows. For a lower bound take a diagonal orthogonal sublattice $K$ of index $b$ in $L$, and place all divisors of $2b\det B_K$ that are prime in $S$. Among its regular children take those fixed by the sign change in its last coordinate. At a prime $p\notin S$ take all primitive projective normals supported on the first $r=n-1$ coordinates and nonisotropic modulo $p$. They give distinct raw kernels in $K_p$ by the unit-coordinate chart. Since $r$ is odd, their number at $p^e$ is exactly $p^{(r-1)e}$, by the finite-field calculation above. For every $t$ prime to $S$, there are therefore at least $t^{n-2}$ such children of index $bt$ in $L$. A rational isometry between two of them stabilizes their restored lattice $K$. Dividing by $|\Gamma_O(K)|$ bounds the number of represented ordinary classes from below. Summing over $t\le X/b$ prime to $S$ gives the lower bound by finite inclusion--exclusion, and the subseries at $n-1$ dominates a positive multiple of the harmonic series on these $t$.
